% Part: first-order-logic
% Chapter: natural-deduction
% Section: identity

\documentclass[../../../include/open-logic-section]{subfiles}

\begin{document}

\olfileid{fol}{ntd}{ide}

\olsection{\usetoken{S}{identity}-সহ \usetoken{P}{derivation}}

!!{identity}-সহ !!^{derivation}s-এর জন্য অতিরিক্ত অনুমান-বিধি দরকার।

\begin{defish}
\AxiomC{}
\RightLabel{\Intro{\eq}}
\UnaryInfC{$\eq[t][t]$}
\DisplayProof
\hfill
\begin{tabular}{r}
\AxiomC{$\eq[t_1][t_2]$}
\AxiomC{$!A(t_1)$}
\RightLabel{\Elim{\eq}}
\BinaryInfC{$!A(t_2)$}
\DisplayProof
\\[3ex]
\AxiomC{$\eq[t_1][t_2]$}
\AxiomC{$!A(t_2)$}
\RightLabel{\Elim{\eq}}
\BinaryInfC{$!A(t_1)$}
\DisplayProof
\end{tabular}
\end{defish}

উপরের বিধিগুলিতে $t$, $t_1$ ও $t_2$ বদ্ধ পদ। \Intro{\eq} বিধি
কোনো অনুমিতি ছাড়াই সরাসরি $\eq[t][t]$ রূপের যেকোনো অভিন্নতার
বক্তব্য !!{derive} করতে দেয়।

\begin{ex}
$s$ ও $t$ বদ্ধ পদ হলে $!A(s), \eq[s][t] \Proves !A(t)$:
\begin{prooftree}
\AxiomC{$\eq[s][t]$}
\AxiomC{$!A(s)$}
\RightLabel{$\Elim{\eq}$}
\BinaryInfC{$!A(t)$}
\end{prooftree}
এটি ``অভিন্নগুলির প্রতিস্থাপনযোগ্যতার নীতি'' বা লাইবনিৎসের সূত্র নামে
পরিচিত হতে পারে।
\end{ex}

\begin{prob}
প্রমাণ করো যে $=$ প্রতিসম ও পরিযায়ী—অর্থাৎ
$\lforall[x][\lforall[y][(\eq[x][y] \lif
    \eq[y][x])]]$ এবং $\lforall[x][\lforall[y][\lforall[z]((\eq[x][y]
    \land \eq[y][z]) \lif \eq[x][z])]]$-এর !!{derivation}s দাও।
\end{prob}

\begin{ex}
আমরা !!{sentence}-টি !!{derive} করি:
\begin{align*}
& \lforall[x][\lforall[y][((!A(x) \land !A(y)) \lif \eq[x][y])]]
\intertext{নিচের !!{sentence} থেকে}
& \lexists[x][\lforall[y][(!A(y) \lif \eq[y][x])]]
\end{align*}
!!{derivation}-টি পিছন দিক থেকে গড়ি:
\begin{prooftree}
\AxiomC{$\lexists[x][\lforall[y][(!A(y) \lif \eq[y][x])]]
\quad \Discharge{!A(a) \land !A(b)}{1}$}
\DeduceC{$\eq[a][b]$}
\DischargeRule{\Intro{\lif}}{1}
\UnaryInfC{$((!A(a) \land !A(b)) \lif \eq[a][b])$}
\RightLabel{\Intro{\lforall}}
\UnaryInfC{$\lforall[y][((!A(a) \land !A(y)) \lif \eq[a][y])]$}
\RightLabel{\Intro{\lforall}}
\UnaryInfC{$\lforall[x][\lforall[y][((!A(x) \land !A(y)) \lif \eq[x][y])]]$}
\end{prooftree}
এবার মুখ্য অনুমিতিটি ব্যবহার করতে হবে। সেটি একটি অস্তিত্বসূচক
!!{formula} বলে মধ্যবর্তী সিদ্ধান্ত~$\eq[a][b]$ !!{derive} করতে
\Elim{\lexists} ব্যবহার করি।
\begin{prooftree}
\AxiomC{$\lexists[x][\lforall[y][(!A(y) \lif \eq[y][x])]]$}
\AxiomC{$\Discharge{\lforall[y][(!A(y) \lif \eq[y][c]])}{2}$}
\noLine
\UnaryInfC{$\Discharge{!A(a) \land !A(b)}{1}$}
\DeduceC{$\eq[a][b]$}
\DischargeRule{\Elim{\lexists}}{2}
\BinaryInfC{$\eq[a][b]$}
\DischargeRule{\Intro{\lif}}{1}
\UnaryInfC{$((!A(a) \land !A(b)) \lif \eq[a][b])$}
\RightLabel{\Intro{\lforall}}
\UnaryInfC{$\lforall[y][((!A(a) \land !A(y)) \lif \eq[a][y])]$}
\RightLabel{\Intro{\lforall}}
\UnaryInfC{$\lforall[x][\lforall[y][((!A(x) \land !A(y)) \lif \eq[x][y])]]$}
\end{prooftree}
ওপরের ডানদিকের উপ-!!{derivation}-টি তার অনুমিতিগুলি ব্যবহার করে
$\eq[a][c]$ ও $\eq[b][c]$ দেখিয়ে সম্পূর্ণ করতে হবে। এর জন্য দুটি
আলাদা !!{derivation}s দরকার। $\eq[a][c]$-এর !!{derivation}-টি হলো:
\begin{prooftree}
\AxiomC{$\Discharge{\lforall[y][(!A(y) \lif \eq[y][c]])}{2}$}
\RightLabel{\Elim{\lforall}}
\UnaryInfC{$!A(a) \lif \eq[a][c]$}
\AxiomC{$\Discharge{!A(a) \land !A(b)}{1}$}
\RightLabel{\Elim{\land}}
\UnaryInfC{$!A(a)$}
\RightLabel{\Elim{\lif}}
\BinaryInfC{$\eq[a][c]$}
\end{prooftree}
$\eq[a][c]$ ও $\eq[b][c]$ থেকে \Elim{\eq} দিয়ে $\eq[a][b]$
!!{derive} করি।
\end{ex}

\begin{prob}
নিচের !!{formula}s-এর !!{derivation}s দাও:
\begin{enumerate}
\item $\lforall[x][\lforall[y][((\eq[x][y] \land !A(x)) \lif !A(y))]]$
\item $\lexists[x][!A(x)] \land \lforall[y][\lforall[z][((!A(y) \land
    !A(z)) \lif \eq[y][z])]] \lif \lexists[x][(!A(x) \land
  \lforall[y][(!A(y) \lif \eq[y][x])])]$
\end{enumerate}
\end{prob}

\end{document}
